\label{chap:solution}

Para abordar as limitações das arquiteturas CNN tradicionais na classificação de bases de dados forenses de ferimentos por arma de fogo altamente desbalanceadas e não-lineares, este trabalho introduz uma nova arquitetura híbrida: a \textbf{TransferKAN}.

\section{A Arquitetura TransferKAN}
\label{sec:transferkan_arch}

A filosofia central por trás da TransferKAN é combinar a comprovada capacidade de extração de características espaciais das Redes Neurais Convolucionais do estado da arte com o potencial de classificação altamente expressivo e eficiente em parâmetros das Redes de Kolmogorov-Arnold (KAN).

Arquiteturas tradicionais, como a ResNet152~\cite{he2016resnet} utilizada no \textit{benchmark} de 2024~\cite{Lira2025deep}, processam uma imagem por meio de uma série de blocos convolucionais para gerar um vetor de características de alta dimensão. Esse vetor é então tipicamente alimentado em uma cabeça de classificação de Perceptron de Múltiplas Camadas (MLP). Embora o \textit{backbone} ResNet seja excepcionalmente hábil em reconhecer padrões visuais (bordas, texturas, formas), a cabeça MLP frequentemente tem dificuldade em mapear essas características em classes discretas quando as fronteiras de decisão são altamente complexas, ruidosas e não-lineares---como é inerentemente o caso das imagens forenses de ferimentos por arma de fogo.

\subsection{Componentes da Arquitetura}

A arquitetura TransferKAN é composta por três estágios sequenciais distintos:
\begin{enumerate}
    \item \textbf{Extrator de Características Espaciais (Backbone ResNet152):} Utilizamos um modelo ResNet152~\cite{he2016resnet} pré-treinado na ImageNet~\cite{deng2009imagenet}. A camada densa totalmente conectada final (o classificador MLP original) é removida programaticamente. Consequentemente, para uma imagem de entrada padrão, o \textit{backbone} produz um vetor de características achatado de tamanho 2048. Isso aproveita o paradigma do Aprendizado por Transferência, evitando a necessidade de treinar um extrator profundo de características do zero em uma base de dados forense limitada.

    \item \textbf{Regularização (Dropout):} Para conectar o \textit{backbone} ao novo classificador e, ao mesmo tempo, prevenir agressivamente o sobreajuste, uma camada de \textit{Dropout} com probabilidade $p=0{,}5$ é introduzida. Isso força a rede a aprender representações redundantes e robustas, em vez de memorizar as amostras de treinamento.

    \item \textbf{Classificador de Rede de Kolmogorov-Arnold (KAN):} A inovação central reside em acoplar um módulo KAN~\cite{liu2024kan} diretamente ao vetor de características. Em vez de uma multiplicação matricial de camada densa com ativações fixas nos nós, a KAN mapeia o vetor de 2048 dimensões para o número desejado de classes de saída (por exemplo, 2 para Entrada/Saída ou 3 para MLSD) por meio de uma camada oculta intermediária de 128 dimensões. As conexões entre essas camadas são modeladas como B-splines aprendíveis.
\end{enumerate}

O fluxo completo de dados da arquitetura proposta é resumido na Figura~\ref{fig:transferkan_arch}.

\begin{figure}[htpb]
\centering
\begin{tikzpicture}[
    >=Stealth,
    node distance=8mm,
    block/.style={rectangle, rounded corners, draw=black, thick, align=center,
                  minimum height=13mm, minimum width=20mm, fill=blue!8},
    backbone/.style={block, fill=green!10, minimum width=24mm},
    kan/.style={block, fill=orange!12, minimum width=24mm},
    io/.style={block, fill=gray!12, minimum width=16mm},
    every node/.style={font=\small},
]
    \node[io] (input) {Imagem de\\entrada\\$224\times224$};
    \node[backbone, right=of input] (resnet) {Backbone\\ResNet152\\(ImageNet)};
    \node[block, right=of resnet] (drop) {Dropout\\$p=0{,}5$};
    \node[kan, right=of drop] (kanlayer) {Classificador\\KAN\\$[2048,128,N]$};
    \node[io, right=of kanlayer] (out) {Logits de\\classe\\($N{=}2$ ou $3$)};

    \draw[->, thick] (input) -- (resnet);
    \draw[->, thick] (resnet) -- node[above, font=\scriptsize] {2048-d} (drop);
    \draw[->, thick] (drop) -- (kanlayer);
    \draw[->, thick] (kanlayer) -- (out);

    \begin{scope}[on background layer]
        \node[draw=green!50!black, dashed, rounded corners, inner sep=3mm,
              fit=(resnet), label={[font=\scriptsize, green!50!black]below:Congelado na Fase~1}] {};
        \node[draw=orange!70!black, dashed, rounded corners, inner sep=3mm,
              fit=(kanlayer), label={[font=\scriptsize, orange!70!black]below:Cabeça treinável}] {};
    \end{scope}
\end{tikzpicture}
\caption{A arquitetura TransferKAN proposta. Um \textit{backbone} ResNet152 pré-treinado na ImageNet extrai um vetor de características de 2048 dimensões, que é regularizado por uma camada de \textit{dropout} ($p=0{,}5$) e mapeado para as classes de saída por um classificador de Rede de Kolmogorov-Arnold com estrutura $[2048,128,N]$ ($N=2$ para Entrada/Saída, $N=3$ para MLSD).}
\label{fig:transferkan_arch}
\end{figure}

\section{Justificativa da Inovação}
\label{sec:justification}

A transição de uma cabeça de classificação MLP para um classificador KAN não é uma mera substituição estrutural; ela representa uma mudança fundamental na forma como a rede modela a fronteira de decisão final.

Ao colocar as funções de ativação não-lineares nas arestas (splines), a KAN molda dinamicamente suas formas de ativação para melhor se ajustar aos dados durante o treinamento. Isso oferece duas vantagens significativas no contexto forense:
\begin{itemize}
    \item \textbf{Discriminação Superior em Dados Desbalanceados:} As splines permitem que a KAN crie fronteiras não-lineares altamente intrincadas ao redor dos agrupamentos das classes minoritárias (como ferimentos de Saída ou disparos Encostados) sem colapsar na classe majoritária, um ponto comum de falha dos MLPs rígidos.
    \item \textbf{Eficiência de Parâmetros:} A KAN requer significativamente menos parâmetros para alcançar um mapeamento matemático altamente complexo em comparação com um MLP, reduzindo o risco de superparametrização e de sobreajuste à base de dados FDCPUnBGunshotDB, que é limitada.
\end{itemize}

Por meio desse projeto arquitetural estratégico, a TransferKAN busca atuar como um discriminador probabilístico otimizado, adaptado especificamente às nuances da patologia forense.
